%  LaTeX support: latex@mdpi.com 
%  For support, please attach all files needed for compiling as well as the log file, and specify your operating system, LaTeX version, and LaTeX editor.

%=================================================================
\documentclass[jmse,article,submit,pdftex,moreauthors]{Definitions/mdpi} 

%=================================================================
% MDPI internal commands
\firstpage{1} 
\makeatletter 
\setcounter{page}{\@firstpage} 
\makeatother
\pubvolume{1}
\issuenum{1}
\articlenumber{0}
\pubyear{2023}
\copyrightyear{2023}
%\externaleditor{Academic Editor: Firstname Lastname}
\datereceived{14 March 2023} 
\daterevised{17 April 2023} 
\dateaccepted{} 
\datepublished{} 
\hreflink{https://doi.org/} % If needed use \linebreak
%\doinum{}

%=================================================================
% Add packages and commands here. The following packages are loaded in our class file: fontenc, inputenc, calc, indentfirst, fancyhdr, graphicx, epstopdf, lastpage, ifthen, lineno, float, amsmath, setspace, enumitem, mathpazo, booktabs, titlesec, etoolbox, tabto, xcolor, soul, multirow, microtype, tikz, totcount, changepage, attrib, upgreek, cleveref, amsthm, hyphenat, natbib, hyperref, footmisc, url, geometry, newfloat, caption

\graphicspath{{figures/}} % path to folder with figures
\usepackage{subcaption}
\usepackage{xcolor}
\usepackage{gensymb} % for \textdegree
\usepackage{tabularx}
\usepackage{makecell}
\usepackage{array}
\usepackage{siunitx}
\usepackage{amsmath}

\usepackage{listings}
\usepackage{xcolor}
\definecolor{deepblue}{rgb}{0,0,0.5}
\definecolor{deepred}{rgb}{0.6,0,0}
\definecolor{deepmagenta}{RGB}{195,12,214}
\definecolor{deepgreen}{RGB}{29,122,30}
\definecolor{mintgreen}{RGB}{192,249,192}
\definecolor{codepurple}{RGB}{204, 0, 153}
\definecolor{LightCyan}{rgb}{0.88,1,1}
\definecolor{GainsboroPurple}{RGB}{247,176,247}
\definecolor{LightSlateGrey}{RGB}{124,118,118}
%    
\lstdefinestyle{mystyle}{
    commentstyle=\color{LightSlateGrey},
    keywordstyle=\color{deepgreen}\bfseries,
    emphstyle=\color{deepred},
    numberstyle=\tiny\color{deepmagenta},
    stringstyle=\color{codepurple},
    basicstyle=\ttfamily\footnotesize,
    breakatwhitespace=false,         
    breaklines=true,                 
    captionpos=b,                    
    keepspaces=true,                 
    numbers=left,                    
    numbersep=5pt,                  
    showspaces=false,                
    showstringspaces=false,
    showtabs=false,                  
    tabsize=2
}
\lstset{style=mystyle}

%=================================================================
% Full title of the paper (Capitalized)
\Title{Computing Vegetation Indices from the Satellite Images Using GRASS GIS Scripts for Monitoring Mangrove Forests in the Coastal Landscapes of Niger Delta, Nigeria}

% MDPI internal command: Title for citation in the left column along the Southwestern Coast of Nigeria
\TitleCitation{Computing Vegetation Indices from the Satellite Images Using GRASS GIS Scripts for Monitoring Mangrove Forests in the Coastal Landscapes of Niger Delta, Nigeria}

% Author Orchid ID: enter ID or remove command
\newcommand{\orcidauthorA}{0000-0002-5759-1089} % Polina Add \orcidA{} behind the author's name
\newcommand{\orcidauthorB}{0000-0002-6461-1551} % Olivier Add \orcidB{} behind the author's name

% Authors, for the paper (add full first names)
\Author{Polina Lemenkova $^{1}$*\orcidA{} and Olivier Debeir $^{1}$\orcidB{}}

%\longauthorlist{yes}

% MDPI internal command: Authors, for metadata in PDF
\AuthorNames{Polina Lemenkova and Olivier Debeir}

% MDPI internal command: Authors, for citation in the left column
\AuthorCitation{Lemenkova, P.; Debeir, O.}
% If this is a Chicago style journal: Lastname, Firstname, Firstname Lastname, and Firstname Lastname.

% Affiliations / Addresses (Add [1] after \address if there is only one affiliation.)
\address{%
$^{1}$ \quad Laboratory of Image Synthesis and Analysis (LISA), École polytechnique de Bruxelles (Brussels Faculty of Engineering), Université Libre de Bruxelles (ULB). Building L, Campus du Solbosch, ULB -- LISA CP165/57, Avenue Franklin D. Roosevelt 50, Brussels 1050, Belgium.
}

% Contact information of the corresponding author
\corres{Correspondence: polina.lemenkova@ulb.be; Tel.: +32 471 86 04 59 (P.L.)}

% Current address and/or shared authorship
%\firstnote{Current address: Affiliation 3.} 

% Abstract (Do not insert blank lines, i.e. \\) 
\abstract{
This paper addresses the issue of the satellite image processing using GRASS GIS in the mangrove forests of the Niger River Delta, southern Nigeria. The estuary of the Niger Delta in the Gulf of Guinea is an essential hotspot of biodiversity in western coast of Africa. At the same time, climate issues and anthropogenic factors pose affect vulnerable coastal ecosystems and result in rapid decline of mangrove habitats. This motivates monitoring of the vegetation patterns by the advanced cartographic methods and data analysis. As a response to this need, this study aims at calculating and mapping several vegetation indices (VI) using scripts as advanced programming methods integrated in geospatial studies. The data include four Landsat 8-9 OLI/TIRS images covering the western segment of the Niger Delta in the Bight of Benin for 2013, 2015, 2021 and 2022. The techniques include the \emph{i.vi}, \emph{i.landsat.toar} and other modules of the GRASS GIS. Based on the GRASS GIS \emph{i.vi} module, ten VI are computed and mapped for the western segment of the Niger Delta estuary: Atmospherically Resistant Vegetation Index (ARVI), Green Atmospherically Resistant Vegetation Index (GARI), Green Vegetation Index (GVI), Difference Vegetation Index (DVI), Perpendicular Vegetation Index (PVI), Global Environmental Monitoring Index (GEMI), Normalized Difference Water Index (NDWI), Second Modified Soil Adjusted Vegetation Index (MSAVI2), Infrared Percentage Vegetation Index (IPVI) and Enhanced Vegetation Index (EVI). The results showed variations in the vegetation patterns in mangrove habitats situated in the Niger Delta over the last decade as well as the increase in urban area (Onitsha, Sapele, Warri and Benin City) and settlements in the Delta State due to urbanisation. The advanced techniques of the GRASS GIS of satellite image processing and analysis enabled to identify and visualise changes in vegetation patterns. Technical excellence of the GRASS GIS in image processing and analysis is demonstrated in the scripts used in this study. 
}

% Keywords
\keyword{Image processing; West Africa; remote sensing; Niger River Delta; R language.} 

\PACS{91.10.Da; 91.10.Jf; 91.10.Sp; 91.10.Xa; 96.25.Vt; 91.10.Fc; 95.40.+s; 95.75.Qr; 95.75.Rs; 42.68.Wt}
\MSC{86A30; 86-08; 86A99; 86A04}
\JEL{Y91; Q20; Q24; Q23; Q3; Q01; R11; O44; O13; Q5; Q51; Q55; N57; C6; C61}

%%%%%%%%%%%%%%%%%%%%%%%%%%%%%%%%%%%%%%%%%%
\begin{document}

\section{Introduction}

\subsection{Background}

This paper presents a GRASS GIS based mapping framework for computing and visualization Vegetation Indices (VI). The GRASS GIS provides a principled way to calculate the VI from the bands of the satellite images using special module \emph{i.vi}. We derive experimental performance bounds of the GRASS GIS for remote sensing data analysis, and demonstrate its utility on a challenging task of environmental monitoring of mangroves ecosystems located in the coastal area of the Niger River Delta in southern Nigeria, West Africa.

Mangroves have a high environmental and socioeconomic value. They provide important functions as habitat for species that live in the intertidal environment \cite{James} and services for human well-being including economic, environmental, and social aspects \cite{AKANNI2018387}. At the same time, recently mangrove ecosystems threatened with extinction which is proved by the increasing rate of deforestation \cite{Enaruvbe2016,Numbere2021,Osuji,Fasona}. Identifying the conditions and dynamics of mangrove ecosystems is possible through the comparison of the VI computed from the bands of the satellite images \cite{HATI2022,ALATORRE201698,SATHISH2021102054}. The VI enable the investigation of the spatial-temporal trends in the distribution of mangroves detected when comparing several images.

Remote sensing data and advanced programming methods provide a valuable combination for monitoring mangrove distribution to evaluate the level of degradation \cite{Hu}. This is possible due to the fundamental feature of the remote sensing data of different absorption and reflectance of solar radiation by various land cover types: water, vegetation, urban spaces, deserts, sands, bare soil and rocks. Moreover, spectral reflectance of diverse types of vegetation vary across various channels of satellite images. This enables to detect various patterns of vegetation and discriminate them from other land cover types. The combination of spectral bands of a satellite image enables to calculate and visualise the distribution of vegetation patterns as characteristics of the Earth's landscapes \cite{CHOPADE2023118839,GITAU2023102898,MISHRA2021107486}.

Comparison of the VI computed from the images taken for several consecutive years enables to evaluate the vegetation patterns \cite{OCHEGE2017211,Lemenkova202229}. This approach builds on quantifying the difference between the values of pixels on bands in a satellite image. Since spectral reflectance of pixels corresponding to vegetation varies in selected channels, e.g., Near-Infrared (NIR) and Red, such properties can be used for discriminating healthy green plants. For instance, the Normalized Difference Vegetation Index (NDVI) is a well known indicator of vegetation health and is widely used for assessment of degradation in sensitive coastal ecosystems through the ratio of NIR and Red \cite{alademomi2020assessing,NSE2020e00599}. Besides the NDVI, other indices are developed which differ in their computational approach, various combinations of Red and NIR bands or addition of other bands, e.g., Green, Blue, SWIR \cite{Adamu,rs10060897}. 

Previous studies used the Geographic Information System (GIS) to quantitatively evaluate and map environmental degradation and decline of vegetation of mangroves \cite{NUMBERE2022433,FASHAE2022e01286,Olorunfemi,Fagbeja,Enaruvbe2015}. However, there is a deficiency in functionality of the traditional GIS: slow-speed data processing, manual operational workflow and subjectivity and biased data analysis. In contrast, using programming techniques support accurate computations and mapping \cite{9883588,Lemenkova202227,9884244,Lemenkova202230,9964623} presenting more advanced tools for environmental monitoring of mangroves. This paper is a continuation of our previous work \cite{technologies11020046} on the use of GRASS GIS scripts for cartographic data processing, where spatial data were processed by the GRASS GIS modules. Herein, we provide an extended application of scripts towards satellite image analysis, and use GRASS GIS for computing the VI to achieve better performance compared to the traditional GIS in terms of both accuracy through automation of data processing, and computational costs due to open source availability of the GRASS GIS. Specifically, we computed several VI to evaluate their performance for monitoring mangrove ecosystems.

\subsection{Study Focus}
Mangroves are a unique halophyte type of the shrubs and forests that have a unique complex root system with prop roots \cite{Numbere}. Mangroves exist only in the saline waters of the intertidal coastal zone of tropics and subtropics. Southern Nigeria, located along the tropical coasts of the Atlantic Ocean, has favourable environmental and climate setting for mangrove distribution. In particular, the Niger River Delta has the most extensive area of mangroves in West Africa \cite{KINAKO197735}, and it is the third largest wetland in the world \cite{ZABBEY2021100571}. A wide variety of mangrove species are adapted to the environmental setting of Nigeria \cite{Ukpong1991,Ufuoma}. For instance, the Niger River Delta (Figure \ref{fig01}) provides a habitat for \emph{Rhizophora} which is the most common species located along the coasts, or shrub genera \emph{Avicennia} and \emph{Laguncularia} that growth in the lagoon areas of southern Nigeria \cite{Akpovwovwo}. The dynamics of mangrove habitats and the distribution of these species depend on various factors: salinity \cite{Amadi,Okafor,Ukpong1997}, distance from coast, repeatability of tidal cycles and freshwater inflow. Other factors include landscape characteristics that vary in the coastal ecosystems of southern Nigeria in the lagoon, mud beach, the Niger River Delta and in the coastal areas that extend along the coasts of the Bight of Benin in the Gulf of Guinea \cite{Ajao}. The mangrove forests have a year-round growing season which continuous vegetation cycle in the mangrove development which includes opening the forest canopy, growing rapidly young trees, and mature phase formed by the large trees \cite{Whitmore}.

\begin{figure}[H]
\centering
	\includegraphics[width=14.0 cm]{Fig_01.jpg}
	\caption{Topographic map of Nigeria with indicated study area. Software: Generic Mapping Tools (GMT) version 6.1.1, \cite{Wessel}. Data source: GEBCO/SRTM \cite{GEBCO,SRTM}. Cartography source: authors.
	\label{fig01}}
\end{figure}

The deforestation of mangrove forests in southern Nigeria has diverse reasons. The most important are the effects from climate change such as temperature and rainfall variability \cite{Adefolalu,AREOLA2018153,w12020385}, because the Niger Delta is the most vulnerable region to climate issues in Nigeria \cite{BALOGUN2022100304,hydrology7010019}. Specifically, the decrease in precipitation and rise in temperatures results in changed balance of the saline-fresh water with regard to hydrology, and variations in the length of periods of dry and wet seasons with regard to soil-landscapes systems. As a consequence, this results in land cover change and deforestation \cite{Akintuyi,KHADIJAT2021983}. Furthermore, environmental processes such as droughts, soil erosion and invasion of external species also affect the health of mangroves ecosystems \cite{UKPONG199761}. Finally, the impacts of wave variability and increased frequency of flash floods changed the shape of the coastlines, contributed to sea level rise and land subsidence, which in turn affected mangrove habitats \cite{DARAMOLA2022102167,OSINOWO2021104471}. 

Another important threat to distribution of mangroves is presented by the anthropogenic activities such as oil and gas exploration in the Niger Delta \cite{su142114256} which results in environmental problems such as oil bunkering, oil spillage \cite{OBIDA2021145854}, refining of crude oil \cite{EYANKWARE2020100293}. Related processes include air pollution due to gas flaring, fossil fuel burning and transportation contribute to the degradation of the coastal landscapes \cite{socsci11110525}. Oil spillage increases the concentration of the polycyclic aromatic hydrocarbons in the sediments \cite{Sojinu}. In turn, the dynamics of the suspended sediments is controlled by different density in fresh and salt waters, energy of waves and tides \cite{Munasinghe}.  In contrast to the turbulent waters, the pollution in the Niger Delta is aggravated by a more complex hydromorphological structure of the tidal network \cite{GEORGE2019103592}. 

Due to a dense pattern of fluvial and tidal channels, bars and inner lakes interspersed in the lower plain of the Niger Delta, polluted water remains stagnant. Furthermore, clayey soil texture in the meander belt of the Niger Delta \cite{Kamalu} creates favourable conditions for the accumulation of pollutants and toxic chemicals \cite{w13223295,toxics11010006,UKHUREBOR2021112872}, and retains retain heavy metals \cite{Amusan,OKOYE2022100222,OKOYE2021111273}. All these factors facilitate the spread of oil into stagnant waters of delta and contribute to the decline of mangroves \cite{ONYIA2020377,ODISU2021111993,ZABBEY2014167}. The habitats of mangroves are also affected by the industrial works such as road construction, housing, clearing waterfront vegetation \cite{Zabbey}, development of seaports which all affected the natural landscapes of the Niger River Delta \cite{OMIUNU198957}. Finally, the use of aquaculture reagents in the Niger Delta lead to the pollution of waters which resulted in decline of mangrove landscapes up to 35\% since 1960s \cite{Godstime}. 

Losses in mangroves in the coastal landscapes at a global scale are reported to be up to 62\% between 2000 and 2016, as a result of the land-use change, expansion of the aquaculture land, fishing and agriculture plantations \cite{Goldberg}. Such activities contribute to the environmental changes in the coastal regions of West Africa \cite{land11111982,su12187349,Babanyara}. Finally, timber production, wood harvesting, expansion of agriculture and bio-fuel plantations, and coastal development contribute to the extinction of mangroves \cite{Feka}. The cumulative effects from all these triggers lead to the decline of mangroves at high rates and reduce the number of their remaining habitats. As a result, mangrove forests are nowadays among the most vulnerable marine ecosystems in West Africa \cite{ORIJEMIE2012360}. 

%%%%%%%%%%%%%%%%%%%%%%%%%%%%%%%%%%%%%%%%%%
\section{Materials and Methods}

The distribution of the mangroves within the coastal landscapes of Nigeria was evaluated by using ten different vegetation indices, namely: Atmospherically Resistant Vegetation Index (ARVI), Green Atmospherically Resistant Vegetation Index (GARI), Green Vegetation Index (GVI), Difference Vegetation Index (DVI), Perpendicular Vegetation Index (PVI), Global Environmental Monitoring Index (GEMI), Normalized Difference Water Index (NDWI), Second Modified Soil Adjusted Vegetation Index (MSAVI2), Infrared Percentage Vegetation Index (IPVI) and Enhanced Vegetation Index (EVI). Technically, the methods of computing the ten VI were implemented using the GRASS GIS scripts. 
%----------- subsection ----------->

\subsection{Data}
Rapid growth of the satellite data derived from the Earth observation sensors enhances the cartographic possibilities of the environmental monitoring of the coastal landscapes. The Landsat 8-9 OLI/TIRS images have become a precious data source for a wide range of geospatial applications, e.g., monitoring the coastal ecosystems that are difficult to access, such as mangrove habitats in southern Nigeria. The advantages of the Landsat scenes consist in multi-spectral characteristics and global temporal and spatial coverage which enable to derive information on the Earth's landscapes. The Landsat 8-9 OLI/TIRS images have 11 spectral bands including visible, thermal and panchromatic ones, which can be used for the analysis of the reflectance spectra for a particular land cover types and their spatial and spectral characteristics.

Specifically for inaccessible areas of mangrove ecosystems in southern Nigeria, the Landsat data provide a precious source of spatial information which enables accurate mapping and analysis of the vegetation patterns. Processing the satellite images by the advanced scripting technologies is an effective and profitable means of data analysis and information retrieval. Due to such reasons, we used the Landsat 8-9 OLI/TIRS satellite images to monitor the decline in mangrove habitats based on the computed and visualised ten VI which cover the western segment of the Niger River Delta. The images were collected for the following years: 2013, 2015, 2021 and 2022, Figure \ref{fig02}. 

The data capture was performed using the EarthExplorer repository of the USGS (\href{https://earthexplorer.usgs.gov/}{https://earthexplorer.usgs.gov/}), with main metadata summarised in Table \ref{tab1}. 

\begin{figure}[H]
\makebox[0.99\linewidth][r]{%
	\begin{subfigure}[b]{.5\textwidth}
		\centering
			\hspace{40pt}\includegraphics[width=6.5cm]{Fig_02a.jpg}
		\caption{2013}
	\end{subfigure}%
	\begin{subfigure}[b]{.5\textwidth}
		\centering
			\hspace{40pt}\includegraphics[width=6.5cm]{Fig_02b.jpg}
		\caption{2015}
	\end{subfigure}%
}\\
\vfill \vspace{1mm}
\makebox[0.99\linewidth][r]{%
	\begin{subfigure}[b]{.5\textwidth}
		\centering
			\hspace{40pt}\includegraphics[width=6.5cm]{Fig_02c.jpg}
		\caption{2021}
	\end{subfigure}
	\begin{subfigure}[b]{.5\textwidth}
		\centering
			\hspace{40pt}\includegraphics[width=6.5cm]{Fig_02d.jpg}
		\caption{2022}
	\end{subfigure}%
}\caption{Landsat 8-9 OLI/TIRS images of Niger Delta in RGB colors for four years: (\textbf{a}) 2013/12/20, (\textbf{b}) 2015/12/26, (\textbf{c}) 2021/12/18, (\textbf{d}) 2022/12/29}\label{fig02}
\end{figure}

Among all the existing Landsat products, the Landsat OLI/TIRS sensor has improved technical and spectral characteristics: narrower spectral bands including Red and NIR channels which are crucial for calculation of the VI. Moreover, it has better calibration, higher radiometric resolution and corrections signal for noise. A thorough comparative analysis of the Landsat OLI/TIRS against earlier sensors is presented in \cite{ROY201657}. Since spectral characteristics of the older Landsat products (TM and MSS) have coarser parameters, we only used the images obtained from the Landsat OLI/TIRS. At the same time, the first Landsat 8 OLI/TIRS sensor was launched on 2013, which explains the choice of the first image selected in this study. Another important parameter for data selection was low cloudiness of the images and comparable span for a short time series. 

The images were imported into the GRASS GIS software and processed for data analysis and visualization using\hl{ }modules \hl{for} satellite images and raster data processing.\hl{ }Mapping mangroves of Nigeria using satellite images and GIS techniques has been attempted and reported earlier \cite{Ozigis,GUNDLACH2022100831,Sobande,ijerph2006030011} with proven effectiveness of the satellite data for \hl{monitoring mangroves}. We continue these and similar studies by applying the programming methods of the GRASS GIS scripts instead of the traditional approaches of GIS for computing the VI of the\hl{ }landscapes \hl{in southern Nigeria}. 

While the Landsat 8-9 OLI/TIRS images formed the dataset for computing and visualising the VI, the GEBCO/SRTM raster data were used for visualising the base map for this study \hl{using the techniques of the Generic Mapping Tools (GMT) scripts} \cite{Lemenkova202228,Lemenkova202301}. To determine the variation in \hl{a} spatial extent of the mangroves in\hl{ }western segment of the Niger Delta, ten VI representing the distribution of healthy vegetation were computed to discriminate \hl{the }swamp forests from the other land cover types \hl{such as }water, urban areas \hl{and} agricultural lands. \hl{The} image processing \hl{was implemented} using the GRASS GIS, version 8.2.1, \href{https://grass.osgeo.org/}{https://grass.osgeo.org/}. The GRASS GIS scripts were used in the \emph{'i.vi'} module which was applied for \hl{computing the VI} along with auxiliary modules. 

\begin{table}[H]
\small
\caption{Satellite images used for computing ten VIs: Landsat 8-9 OLI/TIRS from the USGS. \textsuperscript{1}.\label{tab1}}
	\begin{adjustwidth}{-\extralength}{0cm}
%		\newcolumntype{C}{>{\centering\arraybackslash}X}
		\begin{tabularx}{\fulllength}{ p{1.6cm} | p{1.8cm} | p{7.0cm} | p{4.0cm} | p{2.0cm} }
			\toprule
			\textbf{Date} & \textbf{Spacecraft} & \textbf{Landsat Product ID} & \textbf{Scene ID} & \textbf{Cloudiness} \\
			\midrule
			2013/12/20 & Landsat 8 & \makecell{LC08\_L1TP\_189056\_20131220\_20200912\_02\_T1} & LC81890562013354LGN01 & 6.46 \\
			2015/12/26 & Landsat 8 & \makecell{LC08\_L1TP\_189056\_20151226\_20200908\_02\_T1} & LC81890562015360LGN01 & 17.54 \\
			2021/12/18 & Landsat 8 & \makecell{LC09\_L1TP\_189056\_20211218\_20220121\_02\_T1} & LC91890562021352LGN01 & 32.32 \\
			2022/12/29 & Landsat 9 & \makecell{LC08\_L1TP\_189056\_20221229\_20230104\_02\_T1} & LC81890562022363LGN00 & 0.59 \\
			\bottomrule
		\end{tabularx}
	\end{adjustwidth}
	\noindent{\footnotesize{\textsuperscript{1} The Sensor ID is common for all the scenes: Landsat 8-9 OLI/TIRS (Operational Land Imager and Thermal Infrared Sensor), Collection 2 Level-2. Path/row parameters \hl{are }common for all the images: 189/56. Coverage: Niger Delta, southern Nigeria, \hl{the }Bight of Benin. Image source: \hl{United States Geological Survey} (USGS).}}
\end{table}
%----------- subsection ----------->
\subsection{Data preprocessing}

\hl{The }GRASS GIS code used for data preprocessing is presented in Listing \ref{lst01}. It \hl{includes} the steps of importing Landsat channels by the Geospatial Data Abstraction Library (GDAL) and \hl{selected }auxiliary modules. The GDAL library was applied to import raster images of \hl{the }subset with 11 Landsat bands into the GRASS GIS mapset using \hl{the} \emph{r.in.gdal} module. After \hl{importing}, the data were checked and listed using the \emph{g.list} module which lists the available files in \hl{a} working folder. Then, the metadata of the selected random bands were \hl{inspected using} the \emph{gdalinfo} \hl{module}. The raster metadata were displayed using the \emph{r.info} module for further display of the raster maps. 

The principal approach here is the \hl{conversion} of the digital numbers (DN) of \hl{the }pixel values \hl{into the} reflectance values which were used for computing the VI\hl{, as showed in }Listing \ref{lst01}. Afterwards, the\hl{ }11 Landsat bands \hl{were reformatted and copied }to match the input structure of the \emph{i.landsat.toar} through the \emph{g.copy} module. \hl{The \emph{i.landsat.toar} module was used to remove} the haze effect \hl{from} the original scenes by correction \hl{of }the top-of-atmosphere radiance \hl{in} the \hl{Landsat }bands. The \hl{conversion} of the DN pixel values \hl{into the} reflectance values was \hl{implemented by} using \hl{the }DOS1 method. \hl{Upon} the completing of the data preprocessing, the ten VI were computed using the processed bands by the \emph{i.vi} module, \hl{as }described \hl{below }in relevant subsections. 

%------------------------->
\subsection{Atmospherically Resistant Vegetation Index (ARVI)}

The Atmospherically Resistant Vegetation Index (ARVI) is based on using \hl{the spectral }bands with wavelengths \hl{ranging }between 0.650~{\textmu}m and 0.865~{\textmu}m\hl{. In the Landsat 8-9 OLI sensor, this }corresponds to the Bands Red (0.64-0.67~{\textmu}m) and NIR (0.85-0.88 ~{\textmu}m)\hl{. }Additionally, the Blue band with\hl{ }wavelength of 0.45-0.51~{\textmu}m in the Landsat OLI/TIRS is distracted from the Red \hl{band }according to the\hl{ Equation }\ref{eq1}, \cite{134076,HUETE1997440}:

\begin{equation}\label{eq1}
	ARVI = \frac{NIR - (2.0*Red - Blue)}{NIR + (2.0*Red - Blue)}
\end{equation}

In the GRASS GIS, the ARVI was computed using\hl{ }Listing \ref{lst02}\hl{ presented in the Appendix.} 

The advantages of ARVI is that it reduces the role of \hl{the }atmospheric scattering that affects the images. Such effects are caused by the aerosols of both natural origin (\hl{e.g., }rain and fog) or urban air pollution (\hl{e.g., }dust or smog). Initially proposed and developed for \hl{the }assessment of vegetation from the Earth Observing System (EOS) Moderate Resolution Imaging Spectroradiometer (MODIS) sensor, ARVI can \hl{also }be applied for the Landsat OLI/TIRS sensors since it contains necessary \hl{spectral }bands\hl{: }Blue, Red and NIR\hl{.}
%------------------------->
\subsection{Green Atmospherically Resistant Vegetation Index (GARI)}

The Green Atmospherically Resistant Index (GARI) is more perceptive to a wide range of chlorophyll content in leaves of mangroves\hl{. This }correlates with the difference between the reflectance of Green, Red and NIR bands, \hl{since} the Green band is used in the formula \hl{of GARI, }in addition to Red and NIR \hl{spectral bands. Adding the} Blue \hl{spectral }band makes it less dependent on possible atmospheric effects such as aerosol conditions. The \hl{theoretical computation} of GARI is based on the \hl{Equation} \ref{eq2} developed \hl{initially }by \cite{GITELSON1996289}:

\begin{equation}\label{eq2}
	GARI = \frac{NIR - (Green - (Blue - Red))}{NIR + (Green - (Blue - Red))}
\end{equation} 

The computation of GARI in the GRASS GIS is done according to the embedded \hl{formula} and\hl{ }performed by Listing \ref{lst03}\hl{, as showed in the Appendix.} 

%------------------------->
\subsection{Green Vegetation Index (GVI)}

The technique of the Tasselled Cap transformation in remote sensing \hl{is} originally developed by \cite{KauthThomas}. The Tasselled Cap GVI provides more detailed information on the distinct features \hl{of the }landscapes such as plant canopy, water bodies and bare land or \hl{sandy areas}. The method of \hl{the }Tasselled Cap performs a transformation of the original pixels in bands into a new 4D dimension \hl{where the following } parameters \hl{are separated}: soil brightness index, green vegetation index (GVI), yellow stuff index, and a \hl{those} related to atmospheric effects \cite{Huang2002}. Of these, we used the GVI which is based on the estimation of greenness in leaves of mangrove canopies \hl{as information on values of the weighted sums of pixels derived from the colour composites of the Landsat bands}. The computation of \hl{the }Tasselled Cap Green Vegetation Index (GVI) is based on the \hl{Equation }\ref{eq3} developed by \cite{4157507}:

\begin{equation}\label{eq3}
\begin{split}
	GVI = -0.2848 * Blue - 0.2435 * Green - 0.5436 * Red + \\
	0.7243 * NIR + 0.0840 * SWIR1- 0.1800 * SWIR2
\end{split}
\end{equation}

We \hl{used} Listing \ref{lst04} for the computation of \hl{of the }GVI\hl{ as showed in the Appendix.}

%------------------------->
\subsection{Difference Vegetation Index (DVI)}

The advantage of the Difference Vegetation Index (DVI) \hl{consists in the discrimination of} soil \hl{from} the vegetation biomass. However, the disadvantage is that it does not considers the difference between the spectral reflectance and radiance of soil surface caused by \hl{the }atmospheric effects \cite{TUCKER1979127}. The use of \hl{the }DVI is based on the ratio between the NIR and Red \hl{spectral bands} \cite{Naji2018} \hl{as showed in the Equation} \ref{eq4}:

\begin{equation}\label{eq4}
	DVI = (NIR - Red)
\end{equation}

This is the original formula used later in a modified version for computing the NDVI which is based on the DVI in the adjusted way and applied in many works \cite{TUCKER1979237,RONDEAUX199695,DATT1998111} \hl{as showed in Equation} \ref{eq5}:

\begin{equation}\label{eq5}
	NDVI = \frac{(NIR - Red)}{(NIR + Red)}
\end{equation}

The above formulae are embedded \hl{in the GRASS GIS }in the function 'dvi' \hl{of the module \emph{i.vi} }used \hl{for} computing the DVI\hl{ as illustrated in} Listing \ref{lst05}\hl{.}

%------------------------->
\subsection{Perpendicular Vegetation Index (PVI)}

The use of \hl{the }PVI is based on the formula originally developed by \cite{Richardson}\hl{. }Sensitive to \hl{the }atmospheric variations, \hl{the PVI is intended }to monitor the productivity of range, forest\hl{ }and croplands. Similar to\hl{ }DVI, it\hl{ }measures the perpendicular distance of the pixels from the soil line\hl{, since it is }based on the assumption that soil reflectance supplies the background signal of \hl{the }vegetated surfaces. \hl{The formula used for calculation of the PVI is presented in Equation} \ref{eq6}:

\begin{equation}\label{eq6}
	PVI=\frac{1}{\sqrt{a^2+1}} * (NIR - aR - b),
\end{equation} 

where \emph{a} and \emph{b} are\hl{ }constants \hl{indicating} slope and gradient of the soil line \hl{and \emph{R} is a Red band (for the Landsat OLI/TIRS sensor, it is a spectral Band 4)}. However, the embedded formula of GRASS GIS is slightly modified, because the calculation of the PVI here is made using the concept of "isovegetation". The isovegetation line shows the distribution of area covered by equal types of vegetation. Therefore, the pixels in this case are parallel to the soil line which is assigned to a=1. Hence, the constants for soil are removed and the computation of the PVI is adjusted. In this case, it includes the sine and cosine of the NIR and Red channels using the following formula \hl{showed in the Equation} \ref{eq7}, \cite{PERRY1984169}:

\begin{equation}\label{eq7}
	PVI = sin(a)NIR-cos(a)Red
\end{equation} 

Using the syntax of GRASS GIS, we define the PVI as follows in Listing \ref{lst06}\hl{ shown in the Appendix.}

\subsection{Global Environmental Monitoring Index (GEMI)}

The Global Environmental Monitoring Index (GEMI) \hl{was originally developed by Pinty and Verstraete in 1992} \cite{PintyVerstraete}. \hl{The formula of GEMI }has a nonlinear nature and \hl{is} more complex\hl{, }compared to the previous indices, \hl{Equation }\ref{eq8}. It is designed for global environmental monitoring \hl{using the }remote sensing data, which explains its \hl{improved performance that better }considers all types of vegetations. 

\begin{align}
GEMI = \Bigl( \frac{2*((NIR * NIR)-(Red * Red)) + 1.5*NIR+0.5*Red}{NIR + Red + 0.5} * \nonumber\\
        (1 - 0.25 * \frac{2*[(NIR * NIR)-(Red * Red) + 1.5*NIR+0.5*Red]}{NIR + Red + 0.5} \Bigr ) - \frac{Red - 0.125}{1 - Red} \label{eq8}
\end{align}

The algorithm for computing and plotting the GEMI is presented in Listing \ref{lst07}\hl{ shown in Appendix.}

Although GEMI is in general comparable to NDVI, it shows less sensitivity to atmospheric effects. Instead, it is influenced by the brightness of the bare soil. To these reasons, it does not work well in desert arid and semi-arid regions covered by sands, rocks and soil with sparse or no vegetation. For the case of southern Nigeria, however, GEMI works well since tropical dense forests are suitable for the performance of this index. Moreover, the decline of mangroves is well depicted using comparison of GEMI several satellite images.

\subsection{Normalized Difference Water Index (NDWI)}

The values of the Normalized Difference Water Index (NDWI) are strongly dependent on the water content in the landscapes. Hence, the NDWI serves as an accurate indicator of the water areas and water bodies (rivers, lakes, estuaries, deltas). The NDWI is derived from the NIR and Green bands using the original formula developed by \cite{McFeeters} as a ratio of Green and NIR \hl{spectral bands, as seen in the Equation }{eq9}: 

\begin{equation}\label{eq9}
	NDWI = \frac{green - NIR}{green + NIR}
\end{equation} 

Another formula of the NDWI is proposed by \cite{GAO1996257}, which is a different index yet having the same name, since both indices are developed simultaneously in 1996. The NDWI by \cite{GAO1996257} is focused on identification of water content in the plant leaves \hl{using the formula presented in Equation }{eq10}: 

\begin{equation}\label{eq10}
	NDWI = \frac{NIR - SWIR}{NIR + SWIR}
\end{equation} 

However, this is more actual for the semi-arid and arid regions covered by savannah. In this study, we used the NDWI proposed by \cite{McFeeters} using Green and NIR bands for detecting areas covered by Niger Delta, tributaries and lagoon.  Using the embedded algorithm in the GRASS GIS we can identify the NDWI using its syntax by the script showed in Listing \ref{lst08}\hl{ demonstrated in Appendix.}

\subsection{Second Modified Soil Adjusted Vegetation Index (MSAVI2)}

The Second Modified Soil Adjusted Vegetation Index (MSAVI2) is adjusted to provide accurate data in sparse vegetation or plant canopies with deficit of chlorophyll. \hl{Originally developed as MSAVI by Qi et al. in 1992 }\cite{QI1994119}\hl{, }MSAVI2 \hl{was slightly modified by Qi et al. in 1994 }\cite{Qi1994ExternalFC} \hl{with a novel developed formula presented in Equation }\ref{eq11} \hl{which aims at better identifying the }bare soil between the seedlings during leaf development which is useful for agricultural mapping. 

\begin{equation}\label{eq11}
	MSAVI2 = \frac{1}{2}*\Bigl(2*NIR+1-\sqrt{(2*NIR+1)^2-8*(NIR-red)}\Bigr)
\end{equation} 

\hl{According to the analysis of the parameters of MSAVI (NIR and Red bands in complex combinations) and its documentation, it allows precise detection of the seasonal vegetation patterns since includes a soil factor which varies in moisture by different seasons in Nigeria. Moreover, MSAVI2 works well in the areas of the sparse and scattered vegetation, for instance, occupied by urban infrastructure, covered by buildings and roads in the agglomerations and urbanized areas, where it shows good correlation with the NDVI} \cite{FABIJANCZYK2022100721}. \hl{Therefore, the MSAVI2 is useful for detecting urban growth in the rapidly developing towns of southern Nigeria. In GRASS GIS,} MSAVI2 is \hl{computed} using i.vi module \hl{by} Listing \ref{lst09}.

\subsection{Infrared Percentage Vegetation Index (IPVI)}

The specifics of \hl{the }Infrared Percentage Vegetation Index (IPVI) \hl{consists in the} ratio-based computationally fast approach with values ranging from 0 to 1. As one can see from the formula \hl{in the Equation }\ref{eq12}, the IPVI is based on the \hl{ratio between the Red and NIR spectral bands. With this regard, }it is similar and linearly equivalent to the NDVI, yet \hl{is implemented in a }more straightforward \hl{way, }\cite{CRIPPEN199071}:

\begin{equation}\label{eq12}
	IPVI = \frac{NIR}{(NIR+Red)}
\end{equation} 

To compute and plot the IPVI\hl{, }we \hl{used} the GRASS GIS script presented in Listing \ref{lst10}\hl{ in the Appendix}.

\subsection{Enhanced Vegetation Index (EVI)}

The Enhanced Vegetation Index (EVI) \hl{is} designed to quantify vegetation greenness \hl{by the combination of} NIR, Red and Blue \hl{spectral }bands. \hl{The} extended formula \hl{of EVI }includes the \hl{corrections for the }atmospheric conditions and considers the effects from the canopy background. \hl{Due to such properties,} EVI is well adjusted to \hl{identify the }areas with dense vegetation which is suitable for mangrove mapping \cite{HUETE2002195}.\hl{ }Similar to the NDVI, \hl{the values of EVI }are ranging from -1 to +1, with the \hl{range} indicating healthy vegetation \hl{located} between 0.2 and 0.8. \hl{The formula of EVI is presented in the Equation }\ref{eq13}:

\begin{equation}\label{eq13}
EVI = 2.5 * \frac{NIR - Red}{NIR + 6.0 * Red - 7.5 * Blue + 1.0}
\end{equation} 

\hl{In the GRASS GIS, }EVI was\hl{ technically implemented }using Listing \ref{lst11}.

%----------- section ----------->
\section{Results}
The results of the\hl{} image processing show that the proposed GRASS GIS framework of scripting methods is effective for mapping \hl{the }VI \hl{in western segment of the Niger River Delta. The evaluation of the VI is useful for the} comparative assessment of vegetation \hl{patterns }and environmental monitoring \hl{in coastal Nigeria}. In this section, we evaluate the \hl{results of the implementation} of various VI and \hl{maps} visualised and compared on several years \hl{for} the target region\hl{. }The sidewise \hl{histograms} placed on the left of the legend in each of the image \hl{demonstrate} the data distribution with the \hl{study }area\hl{ for }each index.

%----------- subsection ----------->
\subsection{Difference Vegetation Index (DVI)}

\hl{The DVI is visualised in Figure} \ref{fig03}. 

\begin{figure}[H]
\makebox[1.00\linewidth][r]{%
	\begin{subfigure}[b]{.55\textwidth}
		\centering
			\includegraphics[width=7.3cm]{Fig_03a}
		\caption{2013}
	\end{subfigure}%
	\begin{subfigure}[b]{.55\textwidth}
		\centering
			\includegraphics[width=7.3cm]{Fig_03b}
		\caption{2015}
	\end{subfigure}%
}\\
\vfill \vspace{1mm}
\makebox[1.00\linewidth][r]{%
	\begin{subfigure}[b]{.55\textwidth}
		\centering
			\includegraphics[width=7.3cm]{Fig_03c}
		\caption{2021}
	\end{subfigure}
	\begin{subfigure}[b]{.55\textwidth}
		\centering
			\includegraphics[width=7.3cm]{Fig_03d}
		\caption{2022}
	\end{subfigure}%
}\caption{DVI from the Landsat 8-9 OLI/TIRS of Niger Delta: (\textbf{a}) 2013, (\textbf{b}) 2015 (\textbf{c}) 2021, (\textbf{d}) 2022.}\label{fig03}
\end{figure}

The DVI (Figure \ref{fig03}) is a slope-based type of VI which is more straightforward to compute compared to the NDVI, \hl{which is beneficial for mangroves since they are located at the interface between land and sea}. The DVI is sensitive to the total content of vegetation and density of canopy\hl{. }Moreover, it well distinguishes \hl{the areas }between the\hl{ bare }soil and vegetation \hl{canopies }in general and identifies urban areas \hl{in the small settlements}.

The majority of \hl{the }values \hl{obtained }for 2013 are located within the range of 0.06 to 0.18\hl{ since it is the index sensitive to variations in chlorophyll content}. The interpretation is as followed. \hl{The lowest }values around zero indicate bare soil\hl{; slightly higher} values \hl{0.10 to 0.12 }stand for vegetation canopies, while \hl{negative }values \hl{coloured by blue }show water bodies. The values of 2015  show similar data range with a slight increase of upper values towards 0.20. The distribution of mangrove swamps is visible as large \hl{slate-}blue-coloured spot in the left low corner of the image. The comparison of several images of 2013 and 2015 to 2021 and 2022 shows the decrease in this areas with mangroves replaced by the values dominated by bare soils \hl{and }other types of vegetation\hl{, }coloured by the inclusions of green\hl{. The bright red colours indicate tropical forests located in the northern part of the area.}

%----------- subsection ----------->
\subsection{Atmospherically Resistant Vegetation Index (ARVI)}

The ARVI (Figure \ref{fig04}) is designed for the regions of high atmospheric aerosol content \hl{as a robust modification of the NDVI. }

\begin{figure}[H]
\makebox[1.00\linewidth][r]{%
	\begin{subfigure}[b]{.55\textwidth}
		\centering
			\includegraphics[width=7.5cm]{Fig_04a}
		\caption{2013}
	\end{subfigure}%
	\begin{subfigure}[b]{.55\textwidth}
		\centering
			\includegraphics[width=7.5cm]{Fig_04b}
		\caption{2015}
	\end{subfigure}%
}\\
\vfill \vspace{1mm}
\makebox[1.00\linewidth][r]{%
	\begin{subfigure}[b]{.55\textwidth}
		\centering
			\includegraphics[width=7.5cm]{Fig_04c}
		\caption{2021}
	\end{subfigure}
	\begin{subfigure}[b]{.55\textwidth}
		\centering
			\includegraphics[width=7.5cm]{Fig_04d}
		\caption{2022}
	\end{subfigure}%
}\caption{ARVI computed from the Landsat 8-9 OLI/TIRS images of Niger Delta for four years: (\textbf{a}) 2013, (\textbf{b}) 2015 (\textbf{c}) 2021, (\textbf{d}) 2022.}\label{fig04}
\end{figure}

Therefore\hl{, it} corrects the image scenes for the effects from the atmospheric scattering and \hl{air contamination related to dust and smog in the industrial areas where high concentration of atmospheric aerosols is presented due to the oil exploration}. The\hl{ }sideway histograms\hl{ showing distribution of pixels on }the images\hl{ }shows various values ranging for the years 2013, 2015, 2021 and 2022. \hl{Thus, the} histogram for\hl{ }2013 shows a classic bell-shaped data distribution with values ranging from -0.20 to +0.20 distributed equally on the histogram with a prominent mound in the centre and identical tapering to the left and right flings of the histogram. 

The data for 2015 are distributed from -0.10 to +0.60,\hl{ }showing a more oblique data distribution with more data \hl{distributed }within \hl{the }positive interval \hl{which indicates the increased areas covered by green healthy vegetation}. The data for 2021 show variation with the diapason from -0.20 to 0.80 and the data for 2022 -- from 0.00 to 0.85 with residuals in the negative values that indicate water areas. \hl{This indicate the increase in the areas covered by the tropical forests in the northern part of the area, while decreased area of mangroves. }As mentioned above, the advantages of using \hl{ARVI} for Nigeria is that it reduces the impacts of \hl{the }aerosols such as such as rain, fog, dust, smoke, or air pollution \hl{by using the blue band in making atmospheric corrections on the red band}. This is \hl{beneficial} for Nigeria \hl{in the areas of industrial activities such as oil exploration where higher contamination of air is reported. Another important factor is a }high humidity \hl{which is }around 80\% all year throughout for southern Nigeria, \hl{which makes ARVI beneficial specifically in these regions since it makes corrections for atmospheric scattering caused by high humidity and moisture.}

%----------- subsection ----------->
\subsection{Green Atmospherically Resistant Vegetation Index (GARI)}

The results of the computed Green Atmospherically Resistant Vegetation Index (GARI) (Figure \ref{fig05}) rely on the estimated chlorophyll content in plants\hl{ }linked to photosynthetic activity. Therefore, this index demonstrates the higher positive values in the areas with healthy and dense vegetation canopies which is related to an increase in photosynthetic activity of leaves, while negative values indicate water areas \hl{as seen as} a sharp peak in the histogram shape. \hl{For all the four images, the results of Landsat OLI/TIRS images processing with GARI algorithm produce pixel values ranging from -0.50 to +1.0. Based on the GARI results, the peak values in data variations were analysed to facilitate the visual interpretation of the obtained maps, namely the range of values. Specifically }for 2013, the majority of positive values range from 0.15 to 0.45, for 2015 -- from 0.30 to 0.90 with a more gentle gradual shape of histogram; for 2021 -- the histogram has a rather narrow \hl{shape} with values varying from 0.35 to 0.75, and for 2022 -- from 0.20 to 0.80\hl{, }Figure \ref{fig05}\hl{. }Such variations illustrate changes in vegetation coverage and greenness of canopies in mangroves and other vegetation types in the Niger \hl{River} Delta. 

\hl{In 2013, the mangrove forests were dominated by the high GARI values, i.e., between 0.15 to 0.20. while in the 2015, the low-medium values corresponding to the  mangroves are replaced by the higher values (0.80 to 1.00) that indicate the other vegetation types (coloured by red in Figure} \ref{fig05}). \hl{Further, values of histograms increased which indicate the replacement of the mangroves by other vegetation types (e.g., agricultural crop areas) with bright green canopies that correspond to higher values of GARI. The mangrove habitats with moderate values of GARI is presented by the relatively scattered areas in the southern edge of the study area in the Niger River Delta.}

\hl{Such performance of GARI is explained by its high sensitivity to the }nuances of chlorophyll content in plants\hl{ since the computation of GARI} is based on the use of Green, Blue and\hl{ }NIR \hl{spectral }bands. \hl{Similar to ARVI and in contrast to NDVI, }it is less sensitive to \hl{the }atmospheric effects such as aerosol \hl{due to the} rigorous optimization of \hl{its} algorithm\hl{. }Moreover, GARI is sensitive to the vegetation fraction in the canopy as can be seen in the interspersed pattern of diverse colours which represent different values of the index. As a development of present study, GARI can be further applied for estimation of other biophysical properties of mangrove canopy such as Leaf Area Index \hl{(LAI)}, ratio of absorbed photosynthetically active radiation and net primary production for \hl{the }environmental monitoring of the Niger \hl{River }Delta.

\begin{figure}[H]
\makebox[1.00\linewidth][r]{%
	\begin{subfigure}[b]{.55\textwidth}
		\centering
			\includegraphics[width=7.5cm]{Fig_05a}
		\caption{2013}
	\end{subfigure}%
	\begin{subfigure}[b]{.55\textwidth}
		\centering
			\includegraphics[width=7.5cm]{Fig_05b}
		\caption{2015}
	\end{subfigure}%
}\\
\vfill \vspace{1mm}
\makebox[1.00\linewidth][r]{%
	\begin{subfigure}[b]{.55\textwidth}
		\centering
			\includegraphics[width=7.5cm]{Fig_05c}
		\caption{2021}
	\end{subfigure}
	\begin{subfigure}[b]{.55\textwidth}
		\centering
			\includegraphics[width=7.5cm]{Fig_05d}
		\caption{2022}
	\end{subfigure}%
}\caption{GARI computed from the Landsat 8-9 OLI/TIRS images of Niger Delta for four years: (\textbf{a}) 2013, (\textbf{b}) 2015 (\textbf{c}) 2021, (\textbf{d}) 2022.}\label{fig05}
\end{figure}

%----------- subsection ----------->
\subsection{Green Vegetation Index (GVI)}

The GVI is useful for classifying \hl{the }land cover as it recognises the density and vigour plant in green vegetation (bright red areas in Figure \ref{fig06}), detects stressed plants (light blue areas in Figure \ref{fig06}) and accurately identifies wetland areas (dark areas in Figure \ref{fig06}). The values of GVI range from -0.20 to 0.60 for 2013, and then change to the \hl{range between the} -0.20 \hl{and} 0.50 in 2015, which indicates a slight decrease in values corresponding to the green vegetation, while the urban spaces slightly increase. It can also be seen when comparing the area covered by the Benin City for 2013 and 2015. \hl{Further} analysis of the GVI image for 2021 and 2022 shows the range in values varying from -0.20 to +0.50 \hl{for 2021 and from -1.00 to +0.50 for 2022 with the peaks in histogram in the narrow range of -0.10 to +0.10 for 2021 and from -0.5 to 0.15 for 2022}, respectively. \hl{Here the mean mangrove GVI values are around 10, with maximum value of 15 and minimum value of 5.0. }

\hl{In the mangrove forest coverage in 2013, the obtained values of GVI have a range between the 0.2 to 0.8 while for the later periods (2015, 2021 and 2022) they demonstrate lower values which can be seen on the histograms. }Such fluctuations in the values \hl{of GVI }indicate the environmental changes in the study area, and specifically, the decline of the spatial extent of the mangrove swamps. The reasons for the decline of \hl{the }mangroves areas include \hl{the effects from climate changes and the anthropogenic activities such as }logging, \hl{industrial and oil exploration activities,} farming and fishing. \hl{The} effects from oil exploration and related engineering works \hl{such as }dredging contribute to the decline of mangroves \hl{mostly through water contamination. Moreover, }the expansion of \hl{the }urban areas and related air pollution as well as the effects from climate changes also create\hl{ }negative \hl{conditions for mangrove habitats and all }contribute to the gradual extinction of mangroves.

\begin{figure}[H]
\makebox[1.00\linewidth][r]{%
	\begin{subfigure}[b]{.55\textwidth}
		\centering
			\includegraphics[width=7.5cm]{Fig_06a}
		\caption{2013}
	\end{subfigure}%
	\begin{subfigure}[b]{.55\textwidth}
		\centering
			\includegraphics[width=7.5cm]{Fig_06b}
		\caption{2015}
	\end{subfigure}%
}\\
\vfill \vspace{1mm}
\makebox[1.00\linewidth][r]{%
	\begin{subfigure}[b]{.55\textwidth}
		\centering
			\includegraphics[width=7.5cm]{Fig_06c}
		\caption{2021}
	\end{subfigure}
	\begin{subfigure}[b]{.55\textwidth}
		\centering
			\includegraphics[width=7.5cm]{Fig_06d}
		\caption{2022}
	\end{subfigure}%
}\caption{GVI computed from the Landsat 8-9 OLI/TIRS images of Niger Delta for four years: (\textbf{a}) 2013, (\textbf{b}) 2015 (\textbf{c}) 2021, (\textbf{d}) 2022.}\label{fig06}
\end{figure}

%----------- subsection ----------->

\subsection{Perpendicular Vegetation Index (PVI)}

The distance-based PVI index (Figure \ref{fig07}) computes the perpendicular distance of the pixel from the soil line which linearly correlates with the area covered by \hl{the }vegetation canopy. \hl{Therefore, }higher PVI values\hl{ which are }shown as red colours\hl{ }point at the effects from brighter soil background in the agricultural areas in the Edo State \hl{of Nigeria}. According to the sidewise histograms, the majority of values for 2013 roughly range from 0.10 to 0.22, for 2015 -- from 0.06 to 0.18; for 2021 -- from 0.05 to 0.18 with a less distinct peak; for 2022 -- from 0.06 to 0.22. The disadvantages of the PVI consists in the atmospheric effects where it performs less accurate compared to ARVI and GARI. 

\hl{The} trend to broaden \hl{range} of values \hl{is visible in Figure }\ref{fig07} with inclusion of lower values\hl{. Thus, }the green areas of \hl{the }mangrove swamps \hl{have most values} between 0.20 to 0.30 \hl{compared to the} other types of vegetation (rain forests and woodlands) \hl{which cover} the higher values from 0.30 to 0.80. The decline in mangrove forests is visible in the decrease of dark blue-coloured areas in the lower left segment of the \hl{images}. The PVI corrects for the difference between \hl{the }soil background which \hl{however, has a better performance in }arid and semiarid regions. \hl{As for tropical regions such as }Nigeria with its tropical climate and dense swamp forests, \hl{PVI} is less effective and suitable compared to IPVI and MSAVI2\hl{. The latter two indices }better show the difference in \hl{state }state of the mangrove forests. Besides, the PVI\hl{ }is strongly affected by the atmospheric effects which are significant for this region, since southern Nigeria is located in the wet equatorial climate \hl{where} mean annual temperature \hl{is} 26.9°C and high relative humidity \hl{is }above 88\% \hl{for Lagos}. It is also characterised by \hl{a }high cloudiness throughout the country. 

\begin{figure}[H]
\makebox[1.00\linewidth][r]{%
	\begin{subfigure}[b]{.55\textwidth}
		\centering
			\includegraphics[width=7.5cm]{Fig_07a}
		\caption{2013}
	\end{subfigure}%
	\begin{subfigure}[b]{.55\textwidth}
		\centering
			\includegraphics[width=7.5cm]{Fig_07b}
		\caption{2015}
	\end{subfigure}%
}\\
\vfill \vspace{1mm}
\makebox[1.00\linewidth][r]{%
	\begin{subfigure}[b]{.55\textwidth}
		\centering
			\includegraphics[width=7.5cm]{Fig_07c}
		\caption{2021}
	\end{subfigure}
	\begin{subfigure}[b]{.55\textwidth}
		\centering
			\includegraphics[width=7.5cm]{Fig_07d}
		\caption{2022}
	\end{subfigure}%
}\caption{PVI computed from the Landsat 8-9 OLI/TIRS images of Niger Delta for four years: (\textbf{a}) 2013, (\textbf{b}) 2015 (\textbf{c}) 2021, (\textbf{d}) 2022.}\label{fig07}
\end{figure}

%----------- subsection ----------->
\subsection{Global Environmental Monitoring Index (GEMI)}

The visualised GEMI in Figure \ref{fig08} shows the discrimination between the major vegetation classes and other land cover areas \hl{such as} water bodies (dark purple areas of the Bight of Benin and the Niger River with its tributaries),\hl{ }landscapes of the mangrove forests (lilac colour), fresh water swamps (cyan colour) and rain forests near the area of the Benin City (green colour), \hl{as well as }woodlands with tall grass savannah (bright red colour). According to the sidewise histograms, the majority of values for \hl{year }2013 roughly range from 0.40 to 0.60, for 2015 -- from 0.35 to 0.60; for 2021 -- from 0.32 to 0.62; \hl{and }for 2022 -- from 0.30 to 0.70. GEMI corrects the computation for atmospheric and soil-related changes, mostly in the areas where bare soil is presented, which is identified by the dark pixels on the images. Hence, the \hl{computed }GEMI\hl{ }indicates the dynamics in vegetation activity and \hl{the separated areas of }the wetlands and \hl{the }dry lands\hl{. Furthermore, GEMI highlights the} urban areas (dark blue) separated from the natural landscapes (red colour), \hl{and other land cover types such as }water bodies (dark lilac colour), mangroves with fresh water swamps (light purple) and rainforests (cyan colour). The sparsely vegetation areas and agricultural fields are coloured by yellow to orange hues, while the exposed surfaces are coloured green.

\begin{figure}[H]
\makebox[1.00\linewidth][r]{%
	\begin{subfigure}[b]{.55\textwidth}
		\centering
			\includegraphics[width=7.5cm]{Fig_08a}
		\caption{2013}
	\end{subfigure}%
	\begin{subfigure}[b]{.55\textwidth}
		\centering
			\includegraphics[width=7.5cm]{Fig_08b}
		\caption{2015}
	\end{subfigure}%
}\\
\vfill \vspace{1mm}
\makebox[1.00\linewidth][r]{%
	\begin{subfigure}[b]{.55\textwidth}
		\centering
			\includegraphics[width=7.5cm]{Fig_08c}
		\caption{2021}
	\end{subfigure}
	\begin{subfigure}[b]{.55\textwidth}
		\centering
			\includegraphics[width=7.5cm]{Fig_08d}
		\caption{2022}
	\end{subfigure}%
}\caption{GEMI computed from the Landsat 8-9 OLI/TIRS images of Niger Delta for four years: (\textbf{a}) 2013, (\textbf{b}) 2015 (\textbf{c}) 2021, (\textbf{d}) 2022.}\label{fig08}
\end{figure}

\hl{As discussed above, the distribution and health conditions of mangroves in Niger River Delta exhibits a complex characteristics due to the cumulative effects from climate and anthropogenic factors. Besides several species might be presented including the presence of other vegetation species with similar spectral characteristics which complicate the discrimination of mangroves from similar vegetation types. Furthermore, the distribution of mangroves is controlled by variation in salinity. Therefore, as demonstrated above, the distribution of mangrove habitats varies on the Landsat images taken for different years.}

\subsection{Normalized Difference Water Index (NDWI)}

\hl{The} values of the NDWI range from -0.90 to -0.70 in 2013, from -0.90 to -0.50 \hl{according} to the sideway histograms visualised in Figure \ref{fig09}. \hl{As a general trend, this} indicates a decrease in high values \hl{corresponding the green vegetation canopies. Thus, }for 2021 and 2022, the \hl{range} of values lies within \hl{the interval from }-0.85 to -0.45 for 2021, and from -0.80 to -0.40 for 2022. For the latter, one can be an oblique shape of histogram with a peak visibly inclined towards \hl{the }lower values. Such variations in values correspond to the environmental changes in the Niger Delta River and \hl{a} gradual decline in \hl{the distribution of the }mangrove swamps. The identification of vegetation biomass relies on the contrast between \hl{the spectral reflectance of the land cover types corresponding to the }soil and vegetation. Therefore, \hl{the }NIR and Red wavelengths\hl{ }surpass other \hl{spectral bands} of \hl{the }Landsat \hl{images, which is beneficial }for \hl{the }discrimination of \hl{the }green canopies of mangroves.

\begin{figure}[H]
\makebox[1.00\linewidth][r]{%
	\begin{subfigure}[b]{.55\textwidth}
		\centering
			\includegraphics[width=7.5cm]{Fig_09a}
		\caption{2013}
	\end{subfigure}%
	\begin{subfigure}[b]{.55\textwidth}
		\centering
			\includegraphics[width=7.5cm]{Fig_09b}
		\caption{2015}
	\end{subfigure}%
}\\
\vfill \vspace{1mm}
\makebox[1.00\linewidth][r]{%
	\begin{subfigure}[b]{.55\textwidth}
		\centering
			\includegraphics[width=7.5cm]{Fig_09c}
		\caption{2021}
	\end{subfigure}
	\begin{subfigure}[b]{.55\textwidth}
		\centering
			\includegraphics[width=7.5cm]{Fig_09d}
		\caption{2022}
	\end{subfigure}%
}\caption{NDWI computed from the Landsat 8-9 OLI/TIRS images of Niger Delta for four years: (\textbf{a}) 2013, (\textbf{b}) 2015 (\textbf{c}) 2021, (\textbf{d}) 2022.}\label{fig09}
\end{figure}

\subsection{Second Modified Soil Adjusted Vegetation Index (MSAVI2)}

The general \hl{range} of values of MSAVI2 \hl{varies} from -1 to 1, which is similar to the classic NDVI index. More specifically for the case of southern Nigeria,\hl{ }the data \hl{in 2013 }show \hl{that }the majority \hl{of values }range between 0.05 to 0.30 \hl{according to the histogram}; for 2015 -- from 0.05 to 0.40; for 2021 -- from 0.00 to 0.40 with some small fluctuations of the negative data\hl{ }range \hl{ from} -0.15 to 0.00\hl{; in} 2022\hl{ }data mostly \hl{vary} from -0.10 to 0.50\hl{, Figure }\ref{fig10}. \hl{The values of MSAVI2 in the range 0.2--0.4 indicate areas covered with grasslands and shrubs which are situated mostly in the urbanised parts of the settlements such as Onitsha, Sapele, Warri and the Benin City. Therefore, the use of the MSAVI2 index is profitable for highlighting the heterogeneity of soil and vegetation patterns which are visible in such areas. }Mangrove forests \hl{on the maps of MSAVI2 }are coloured by the light orange colour (0.05 to 0.10); water areas -- by \hl{the }dark orange \hl{tones} (0.2 to 0.4); fresh water swamps \hl{--} by light green; rainforests \hl{--} by blue colours (0.20 to 0.25)\hl{, }urban spaces \hl{--} by orange (0.1 to 0.2); and agricultural lands \hl{--} by yellow colours (0.05 to 0.10). The \hl{range} of \hl{the slightly }negative values (-0.20 to -0.10) indicates bare soil and \hl{areas of the }exposed\hl{ }land, Figure \ref{fig10}.

\begin{figure}[H]
\makebox[1.00\linewidth][r]{%
	\begin{subfigure}[b]{.55\textwidth}
		\centering
			\includegraphics[width=7.5cm]{Fig_10a}
		\caption{2013}
	\end{subfigure}%
	\begin{subfigure}[b]{.55\textwidth}
		\centering
			\includegraphics[width=7.5cm]{Fig_10b}
		\caption{2015}
	\end{subfigure}%
}\\
\vfill \vspace{1mm}
\makebox[1.00\linewidth][r]{%
	\begin{subfigure}[b]{.55\textwidth}
		\centering
			\includegraphics[width=7.5cm]{Fig_10c}
		\caption{2021}
	\end{subfigure}
	\begin{subfigure}[b]{.55\textwidth}
		\centering
			\includegraphics[width=7.5cm]{Fig_10d}
		\caption{2022}
	\end{subfigure}%
}\caption{MSAVI2 computed from the Landsat 8-9 OLI/TIRS images of Niger Delta for four years: (\textbf{a}) 2013, (\textbf{b}) 2015 (\textbf{c}) 2021, (\textbf{d}) 2022.}\label{fig10}
\end{figure}

In contrast to SAVI, MSAVI reduces the effects from the bare soil which is especially actual for the agricultural areas. Thus, it presents the improved version of the Soil Adjusted Vegetation Index (SAVI). In turn, SAVI \hl{better }adjusts \hl{the }values in \hl{the }areas with sparse vegetation coverage and low density of \hl{the }plant canopy, which\hl{ }helps to \hl{indicate the} areas with \hl{the }decreased vegetation\hl{ }or high deforestation rates. Since MSAVI2 uses the combination of the ratio of NIR (wavelength 0.85-0.88~{\textmu}m) and Red (wavelength 0.64-0.67~{\textmu}m) where the values reflect the amount of \hl{the }chlorophyll presented in the plant canopy, it is suitable for estimating \hl{the }biomass, grass/plant canopy or \hl{LAI}, because healthy moist vegetation of mangroves reflects high in \hl{the }NIR spectrum while absorbs more in \hl{the }Red. The contrast between these values enables to highlight the areas of \hl{the }distribution of mangrove \hl{habitats}.

\subsection{Infrared Percentage Vegetation Index (IPVI)}

The distinct feature of the IPVI (Figure \ref{fig11}) consists in \hl{its range which is }always positive\hl{. Thus, the data are varying }from 0 to +1, unlike the \hl{typical range} of -1 to +1 for the NDVI and other indices. Specifically, according to the histogram, for the year 2013, the majority of \hl{the IPVI }values\hl{ }range from 0.50 to 0.70; for 2015 -- from 0.50 to 0.80 \hl{with a} more sharp increase in values; \hl{for} 2021 -- from 0.70 to 0.90 and for 2022 - from 0,60 to 0.90. Thus, \hl{one} can \hl{note} the trend of \hl{the generally }increasing higher values that indicate rainforests while the area of mangroves \hl{declines}. The mangroves are depicted \hl{by the} dark blue colours, gradually replaced by the areas of light green colour when comparing the data for 2013 with other years and \hl{specifically, }2022, Figure \ref{fig11}.

\begin{figure}[H]
\makebox[1.00\linewidth][r]{%
	\begin{subfigure}[b]{.55\textwidth}
		\centering
			\includegraphics[width=7.5cm]{Fig_11a}
		\caption{2013}
	\end{subfigure}%
	\begin{subfigure}[b]{.55\textwidth}
		\centering
			\includegraphics[width=7.5cm]{Fig_11b}
		\caption{2015}
	\end{subfigure}%
}\\
\vfill \vspace{1mm}
\makebox[1.00\linewidth][r]{%
	\begin{subfigure}[b]{.55\textwidth}
		\centering
			\includegraphics[width=7.5cm]{Fig_11c}
		\caption{2021}
	\end{subfigure}
	\begin{subfigure}[b]{.55\textwidth}
		\centering
			\includegraphics[width=7.5cm]{Fig_11d}
		\caption{2022}
	\end{subfigure}%
}\caption{IPVI computed from the Landsat 8-9 OLI/TIRS images of Niger Delta for four years: (\textbf{a}) 2013, (\textbf{b}) 2015 (\textbf{c}) 2021, (\textbf{d}) 2022.}\label{fig11}
\end{figure}

Although the IPVI is functionally similar to the NDVI, it is better adjusted for \hl{detecting the} vegetation in small villages and urban spaces due to \hl{a} more fine gradation of \hl{the }positive values. Moreover, the \hl{analysis of values of the }IPVI enable\hl{ }to better identify smaller spots of urban vegetation coverage as well as mangrove wetlands in the Niger Delta. Further,\hl{ }IPVI is well adjusted \hl{to classify} urban ecosystems since the areas of \hl{the }main cities in the study area are distinguishable from the nature vegetation: Warri, Onitsha and Benin City as the largest cities\hl{, and }minor \hl{settlements} such as Agbor, Sapele and Adanisievbe villages in \hl{the }Delta State.

\subsection{Enhanced Vegetation Index (EVI)}

The values of EVI (Figure \ref{fig12}) quantify vegetation greenness in \hl{a general} range of -1 to +1\hl{. }Specifically for the case of Nigeria, we have the most of values distributed between 0.00 to 0.30 in 2013, while for 2015\hl{, }the highest values increase up to 0.40 with remaining 0.00 as the lowest level. The data range for the recent years \hl{varies} from 0.00 to 0.40 \hl{in} 2021 and \hl{from }0.05 to 0.45 \hl{in} 2022. Negative values indicate water areas in all \hl{the }cases while mangrove fields coloured by \hl{the }bright green correspond to values \hl{ranging }from 0.10 \hl{to} 0.15, Figure \ref{fig12}.  

\begin{figure}[H]
\makebox[1.00\linewidth][r]{%
	\begin{subfigure}[b]{.55\textwidth}
		\centering
			\includegraphics[width=7.5cm]{Fig_12a}
		\caption{2013}
	\end{subfigure}%
	\begin{subfigure}[b]{.55\textwidth}
		\centering
			\includegraphics[width=7.5cm]{Fig_12b}
		\caption{2015}
	\end{subfigure}%
}\\
\vfill \vspace{1mm}
\makebox[1.00\linewidth][r]{%
	\begin{subfigure}[b]{.55\textwidth}
		\centering
			\includegraphics[width=7.5cm]{Fig_12c}
		\caption{2021}
	\end{subfigure}
	\begin{subfigure}[b]{.55\textwidth}
		\centering
			\includegraphics[width=7.5cm]{Fig_12d}
		\caption{2022}
	\end{subfigure}%
}\caption{EVI computed from the Landsat 8-9 OLI/TIRS images of Niger Delta for four years: (\textbf{a}) 2013, (\textbf{b}) 2015 (\textbf{c}) 2021, (\textbf{d}) 2022.}\label{fig12}
\end{figure}

\subsection{Statistical Analysis}

\hl{The statistical analysis on the size and number of pixels representing various classes over the study area was performed in R using K-means method applied for Landsat bands 4-3-2 in natural colour composites. 10 vegetation classes corresponding to the coastal landscapes in southern Nigeria, Table }\ref{tab2}. \hl{The results of the classification by K-means clustering show that 10 distinct classes including mangrove classes can be derived from the Landsat OLI/TIRS data. These include the following land cover classes: 1) mangrove swamps; 2) water areas (river and shallow areas); 3) water areas (coastal areas of the Bight of Benin); 4) forests; 5) built-up urban areas and settlements; 6) bare soil and exposed land surfaces; 7) woodland; 8) agricultural areas; 9) farmland/grasslands; 10) other vegetation types. The classification is based on the modified existing studies} \cite{rs12213619,nhess1433172014,Eyoh,Udoka}. \hl{The table is sorted by the ascending values of pixels.} 

\hl{The land cover analysis of the Niger Delta was implemented in R using the K-means clustering method of the unsupervised automatic classification. The approach of the land cover types identification by this method is performed by the partition of the images through sampling the pixels evaluated for their DN values and assigning them to the target land  cover classes. It shows a distribution of pixels by clusters of mangroves and other land cover types using clustering of the four images in Landsat OLI/TIRS band combinations 4-3-2, i.e., natural band composites. The land cover categories corresponding to the mangroves showed gradual losses from 535828.3 in 2013 (ID Class 3 for 2013), to 432318.7 in 2018 (ID Class 4 for 2018), 408108.2 in 2021 (ID Class 5 for 2021) and 416219.5 in 2022 (ID Class 4 for 2022), Table} \ref{tab2}.

\begin{table}[H]
\caption{\hl{Results of the land cover class computations of the Landsat-8 OLI/TIRS satellite images} \textsuperscript{1}.\label{tab2}}
\begin{adjustwidth}{-\extralength}{0cm}
\newcolumntype{C}{>{\centering\small\arraybackslash}X}
    \begin{tabularx}{\fulllength}{| CCC | CC | CC | CC |}
    \toprule
        \bf{Class} & \multicolumn{2}{c|}{\bf{2013}} & \multicolumn{2}{c|}{\bf{2015}} & \multicolumn{2}{c|}{\bf{2021}} & \multicolumn{2}{c|}{\bf{2021}} \\ \hline
        \emph{ID} & \emph{nr. of pixels} & \emph{cluster size} & \emph{nr. of pixels} & \emph{cluster size} & \emph{nr. of pixels} & \emph{cluster size} & \emph{nr. of pixels} & \emph{cluster size} \\ \hline
        \cline{2-9}
        1 & 0.0 & 1 & 99304.0 & 2 & 0.0 & 1 & 0.0 & 1 \\ \hline
        2 & 24202.5 & 2 & 793875.9 & 3 & 179803.3 & 3 & 324638.7 & 1 \\ \hline
        3 & 535828.3 & 4 & 768423.9 & 3 & 249647.7 & 6 & 367419.8 & 3 \\ \hline
        4 & 686438.4 & 8 & 432318.7 & 4 & 305940.8 & 11 & 416219.5 & 6 \\ \hline
        5 & 740729.6 & 9 & 426391.2 & 8 & 408108.2 & 11 & 556665.7 & 9 \\ \hline
        6 & 783803.7 & 9 & 291778.8 & 10 & 603527.3 & 12 & 647066.0 & 9 \\ \hline
        7 & 906337.9 & 10 & 263588.5 & 14 & 628916.8 & 12 & 750900.1 & 12 \\ \hline
        8 & 915597.2 & 16 & 259272.9 & 16 & 629041.5 & 13 & 984741.4 & 13 \\ \hline
        9 & 999878.2 & 18 & 1483831.7 & 20 & 645075.3 & 14 & 1162375.3 & 18 \\ \hline
        10 & 2292524.5 & 23 & 1094319.1 & 20 & 1372956.2 & 17 & 1516837.6 & 18 \\ 
        \bottomrule
    \end{tabularx}
\end{adjustwidth}
\noindent{\footnotesize{\textsuperscript{1} \hl{The statistical summary is computed within cluster sum of squares by cluster using K-means performed in R library 'raster'. The table is sorted by ascending values of the pixels in land cover types}.}}
\end{table}

\hl{The heat map graph shows the confusion matrix plotted using Python showing the probability cases for pixel distribution according to land cover classes (10 selected ID classes) in the Niger Delta, Nigeria, Figure }\ref{fig13}. 

\begin{figure}[H]
\makebox[0.95\linewidth][r]{%
	\begin{subfigure}[b]{.6\textwidth}
		\caption{}\vspace{1mm}
		\centering
			\includegraphics[width=0.95\textwidth]{Fig_13a.png}
	\end{subfigure}%
	\begin{subfigure}[b]{.6\textwidth}
	\caption{}\vspace{1mm}
		\centering
			\includegraphics[width=0.95\textwidth]{Fig_13b.png}
	\end{subfigure}%
}
\caption{\hl{Correlation matrices of land cover classes (10 ID classes) in the Niger River Delta between 2013 and 2022 based on the Landsat 8-9 OLI/TIRS images: (\textbf{a}) Kendall correlation for average values; (\textbf{b}) Spearman correlation for pixels by bands 4 (Red), 3 (Green) and 2 (Blue) used for classification of land cover types. Visualization: Python. Source: authors}.}\label{fig13}
\end{figure}

\hl{The data were evaluated between 2013 and 2022 based on clustering of the Landsat 8-9 OLI/TIRS images. The Kendall correlation shows the average values; while the Spearman correlation shows the pixels by bands 4 (Red), 3 (Green) and 2 (Blue) used for classification of land cover types. The correlation plots showing probability cases for pixel distribution according to land cover classes (10 selected ID classes) in the Niger Delta is shown in  Figure }\ref{fig13}. \hl{It shows the data between 2013 and 2022 based on the results of the clustering of the Landsat 8-9 OLI/TIRS images}. 

\begin{table}[H]\scriptsize
\caption{\hl{Results of the land cover class classification by pixel content in Band 4 (Red), Band 3 (Green) and Band 2 (Blue) of the Landsat-8 OLI/TIRS satellite images}.\label{tab3}}
\begin{adjustwidth}{-\extralength}{0cm}
\newcolumntype{C}{>{\centering\footnotesize\arraybackslash}X}
    \begin{tabularx}{\fulllength}{| C | CCC | CCC | CCC | CCC |}
    \toprule
        \bf{Class} & ~ & \bf{2013} & ~ & ~ & \bf{2018} & ~ & ~ & \bf{2021} & ~ & ~ & \bf{2022} & ~ \\ \hline
         \emph{ID} & \emph{B4} & \emph{B3} & \emph{B2} & \emph{B4} & \emph{B3} & \emph{B2} & \emph{B4} & \emph{B3} & \emph{B2} & \emph{B4} & \emph{B3} & \emph{B2} \\ \hline
        1 & 9362.333 & 9803.222 & 8742.333 & 10352.65 & 10539.35 & 8699.400 & 9150.500 & 9643.167 & 8150.667 & 11409.333 & 11031.333 & 9693.000 \\ \hline
        2 & 9631.111 & 9699.000 & 8448.778 & 10848.50 & 11077.50 & 9607.300 & 11394.000 & 11449.667 & 9541.667 & 9672.556 & 9612.056 & 8631.333 \\ \hline
        3 & 12549.750 & 12197.500 & 10335.750 & 10879.55 & 10833.90 & 8936.950 & 9900.929 & 10211.571 & 8387.929 & 23552.000 & 22211.000 & 18896.000 \\ \hline
        4 & 8920.261 & 9311.826 & 8221.783 & 9816.75 & 10058.75 & 8333.625 & 10262.692 & 10408.692 & 8676.462 & 10117.556 & 9809.111 & 8729.111 \\ \hline
        5 & 10144.625 & 10144.875 & 8946.250 & 9825.25 & 10284.50 & 8776.500 & 9463.333 & 9956.667 & 8546.000 & 8466.667 & 8940.750 & 8148.417 \\ \hline
        6 & 8151.000 & 9824.000 & 8126.500 & 10372.43 & 10680.93 & 9107.500 & 10811.182 & 10739.182 & 8965.091 & 8179.333 & 8705.333 & 7906.833 \\ \hline
        7 & 8658.938 & 9129.000 & 8067.750 & 11872.33 & 11513.67 & 9631.333 & 9408.917 & 9741.917 & 8132.583 & 8887.889 & 9152.389 & 8166.667 \\ \hline
        8 & 8276.889 & 8930.222 & 8045.611 & 10080.38 & 10360.25 & 8400.188 & 15769.000 & 14325.000 & 12446.000 & 9110.308 & 9375.154 & 8483.538 \\ \hline
        9 & 8939.500 & 9494.900 & 8524.100 & 11897.50 & 11900.00 & 10818.000 & 9763.182 & 10158.000 & 8879.636 & 10554.909 & 10198.455 & 9056.000 \\ \hline
        10 & 9217.000 & 10015.000 & 5465.000 & 13052.00 & 12427.00 & 10482.333 & 9559.706 & 9997.765 & 8030.059 & 9374.222 & 9420.889 & 8187.556 \\ \hline
\bottomrule
\end{tabularx}
\end{adjustwidth}
\end{table}

%%%%%%%%%%%%%%%%%%%%%%%%%%%%%%%%%%%%%%%%%%
\section{Discussion}

\hl{This paper demonstrates that scripting techniques of the GRASS GIS used with freely available Landsat OLI/TIRS satellite images are beneficial for evaluating the spatial patterns in the vegetation areas of coastal Nigeria. }The results obtained from the satellite image processing and computing\hl{ }the VI for assessing mangrove distribution shows that MSAVI2 (Figure \ref{fig10}) and GVI (Figure \ref{fig06}) indices ranked the highest with the most accurate presented areas of \hl{the }mangrove fields\hl{, }while EVI (Figure \ref{fig12}) and GEMI (Figure \ref{fig08}) were judged to rank the lowest.\hl{ }Changes in mangrove fields for the four different time periods estimated by the atmospherically adjusted indices ARVI and GARI are presented in Figure \ref{fig04} and Figure \ref{fig05} showing the VI corrected for the atmospheric effects. Here the distribution of the mangrove fields (showed by the light orange colours) recorded the highest value in 2013\hl{, }while the lowest value \hl{was }in 2022 due to the declined area of distribution \hl{of mangrove habitats}, as discussed above. The DVI (Figure \ref{fig03}) and PVI (Figure \ref{fig07}) indices demonstrated the comparable results with visible reduction of mangroves, showed as the reduced area\hl{ }in both the PVI and DVI. Nevertheless, the use of \hl{the }DVI is less effective for mapping mangroves and densely vegetated areas of tropical forests, since this index is better \hl{identifies} urban areas \hl{and settlements}. The NDWI (Figure \ref{fig09}) and the IPVI (Figure \ref{fig11}) also highlighted the decline of mangrove swamps as areas in the lower left segment of the images coloured by yellow green for the NDWI and light blue for the IPVI.

The results of the classified images are also confirmed by the existing studies on changes in mangrove habitats in southern Nigeria reporting the environmental degradation \cite{BAMIDELE2023103274,TOBORE2022158}. For instance, \cite{Godstime2007} confirms the losses in mangroves in southern Nigeria since mid-1980s through early 2000s due to \hl{the urbanisation}, dredging, and oil and gas industrial activities. The effects from \hl{the rapid }urbanisation on deforestation are also reported in other relevant works \cite{AyanladeDrake,su141811289} and reflected in \hl{the }undertaken measures on creating \hl{the }sustainable green infrastructure \cite{ADEGUN2021e01044}. \hl{Despite the }environmental value, mangrove forests along the West African coasts are subject to a range of natural threats \hl{which motivates the environmental monitoring}. The environmental mapping of the coastal landscapes of the Niger River Delta proposed and presented above using \hl{the application of the }GRASS GIS \hl{modules }has provided the capability to integrate \hl{the }remote sensing multi-temporal data covering target region in southern Nigeria\hl{ }with scripting techniques. 

Given the level of mangrove degradation in the Niger Delta, there is a need for mangrove conservation in \hl{the }coastal communities \cite{ONYENA2020e00961}. With this regard, \hl{the presented cartographic} framework \hl{base don the use of the GRASS GIS }contributes to this issue by presenting a case of computed and mapped VI showing recent changes in vegetation communities of \hl{the }southern Nigeria. To evaluate these changes, we visualised ten vegetation indices using \hl{the }remote sensing data by scripting techniques of \hl{selected modules of the }GRASS GIS. To this end\hl{, }we\hl{ }applied and proposed several modules of the GRASS GIS including \hl[{the} 'i.vi' \hl{which is specially adjusted }for satellite image processing\hl{. Using these tools, we performed the }analysis \hl{of the vegetation patterns in the coastal regions of Nigeria by the }Landsat 8-9 OLI/TIRS scenes\hl{. The VI were computed }on the western segment of \hl{the }Niger River Delta \hl{situated along the Bight of Benin in the Gulf of Guinea,} southern Nigeria. Ten\hl{ }computed and visualised \hl{VI were processed }using GRASS GIS scripts\hl{. A }comparative analysis of the environmental changes \hl{was applied for } the \hl{following }vegetation \hl{indices}: ARVI, GARI, GVI, DVI, PVI, GEMI, NDWI, MSAVI2, IPVI and EVI. The use of the advanced scripting methods for mapping changes in the VI within a western segment of the Niger River Delta revealed a declining mangrove swamp habitats in southern Nigeria. 

%%%%%%%%%%%%%%%%%%%%%%%%%%%%%%%%%%%%%%%%%%
\section{Conclusions}

\hl{Identifying spatiotemporal trends in the vegetation patterns is one of the fundamental problems in the environmental assessment. Using} several satellite images \hl{for computing} a time series \hl{of VI }helps to estimate the extent to which\hl{ }mangrove communities \hl{are }declining\hl{. The use of }time series of \hl{the satellite-derived VI }to perform change detection in \hl{the distribution of mangroves }through \hl{the }comparative analysis of \hl{images is reported earlier} \hl{to show} temporal variations in the\hl{ }Niger Delta.

\hl{ In this paper, we illustrated the benefits of usage of scripts for satellite image processing with a special aim of computing and visualization the vegetation indices in mangrove landscapes of Nigeria. Specifically, we applied the GRASS GIS for computing the VI aimed at evaluating the health and greenness in vegetation patterns in southern Nigeria. The use of the GRASS GIS programming scripts is an efficient optimization method of remote sensing data processing. Its benefits for computing the VIs consist in a superior performance with respect to the speed and accuracy of mapping compared to the traditional GIS software. Using modules of the GRASS GIS, we have extracted information from the Landsat OLI/TIRS images for analysis of vegetation patterns in southern Nigeria and compared the results for several years.}

The proposed GRASS GIS framework aimed at image processing and analysis\hl{. Computing the} vegetation \hl{indices} using algorithms embedded in 'i.ivi' and scripting approach \hl{is} integrated consistently for \hl{the }comparative analysis of time series of the Landsat 8-9 OLI/TIRS images covering \hl{the }Niger \hl{River }Delta. \hl{Ten} diverse vegetation indices \hl{were calculated with aim to find the }variations in the distribution of \hl{the vegetation patterns and in particular, }mangrove \hl{habitats of the coastal Nigeria. The performance of the vegetation indices was compared and discussed with regard to suitability of indices to identify mangrove habitats and to discriminate them} against other land cover types \hl{distributed along the coasts of southern Nigeria}. 

To produce \hl{the }colour-consistent images \hl{and evaluate the spatio-temporal changes}, the maps derived from the computed \hl{vegetation indices} were \hl{visualised} in the identical colour ranges \hl{for} each set of the four images. The algorithms of the GRASS GIS demonstrated powerful functionality in satellite image processing \hl{which resulted in the} accurately produced \hl{series of }maps \hl{based on the }calculated VI in \hl{the }western segment of the Niger Delta for years 2013, 2015, 2021 and 2022\hl{, respectively. This allowed for a better understanding of vegetation patterns derived from the analysis of the remote sensing data. Moreover, image processing implemented by the GRASS GIS scripts supports effective environmental monitoring due to rapid and accurate mapping. As a result, this present novel cartographic results as a series of VI-based maps which can be used as a support to protecting sensitive ecosystems in the Niger River Delta. Thus, mapping the VI helps to evaluate the vegetation health as information for adaptive planning and identifying the decline in mangrove habitats which play a crucial role in the environmental sustainability of Nigeria.} 

\hl{Environmental monitoring }of the\hl{ }coastal areas of West Africa supports\hl{ }measures \hl{aimed at} minimising the threats to mangroves and \hl{maintaining the} balance in the human-nature existence for sustainable\hl{ }development in Nigeria. Therefore, \hl{applying and practically evaluating} the programming methods for remote sensing data processing \hl{supports} research on mangrove distribution through \hl{the} advanced cartographic tools. \hl{To} achieve robustness in satellite image analysis, we \hl{used} the GRASS GIS modules based on \hl{the }scripting approach for automation of\hl{ }data processing.

%%%%%%%%%%%%%%%%%%%%%%%%%%%%%%%%%%%%%%%%%%
\authorcontributions{Supervision, conceptualization, methodology, software, resources, funding acquisition, and project administration, O.D.; writing---original draft preparation, methodology, software, data curation, visualization, formal analysis, validation, writing---review and editing, and investigation, P.L. All authors have read and agreed to the published version of the manuscript.}

\funding{The publication was funded by the Editorial Office of JMSE, Multidisciplinary Digital Publishing Institute (MDPI), by providing 100\% discount for the APC of this manuscript. This project was supported by the Federal Public Planning Service Science Policy or Belgian Science Policy Office, Federal Science Policy - BELSPO (B2/202/P2/SEISMOSTORM).}

\institutionalreview{Not applicable.}

\informedconsent{Not applicable.}

\dataavailability{Not applicable} 

\acknowledgments{The authors thank the anonymous reviewers for reading, suggestions and comments that improved an earlier version of this manuscript.}

\conflictsofinterest{The authors declare no conflict of interest.} 

\abbreviations{Abbreviations}{
The following abbreviations are used in this manuscript:\\

\noindent 
\begin{tabular}{@{}ll}
ARVI & Atmospherically Resistant Vegetation Index \\
DN & Digital Number \\
DVI & Difference Vegetation Index \\
EOS & Earth Observing System \\  
EVI & Enhanced Vegetation Index \\
GARI & Green Atmospherically Resistant Vegetation Index \\
GDAL & Geospatial Data Abstraction Library \\
GEMI & Global Environmental Monitoring Index \\
GIS & Geographic Information System \\
GMT & Generic Mapping Tools\\
GRASS GIS & Geographic Resources Analysis Support System GIS \\
GVI & Green Vegetation Index \\
IPVI & Infrared Percentage Vegetation Index \\
Landsat 8-9 OLI/TIRS & Landsat 8-9 Operational Land Imager and Thermal Infrared \\
\hl{LAI} & \hl{Leaf Area Index} \\
MODIS & Moderate Resolution Imaging Spectroradiometer \\
MSAVI2 & Second Modified Soil Adjusted Vegetation Index \\
NDVI & Normalized Difference Vegetation Index \\
NDWI & Normalized Difference Water Index \\
NIR & Near Infrared \\
PVI & Perpendicular Vegetation Index \\
SAVI & Soil Adjusted Vegetation Index \\ 
\end{tabular}
}

\appendixtitles{no} % Leave argument "no" if all appendix headings stay EMPTY (then no dot is printed after "Appendix A"). If the appendix sections contain a heading then change the argument to "yes".
\appendixstart
\appendix

\section[\appendixname~\thesection]{GRASS GIS scripts used for computing and plotting the vegetation indices}

\subsection[\appendixname~\thesubsection]{GRASS GIS script for correction of the top-of-atmosphere radiance}

\begin{lstlisting}[language=bash,caption=GRASS GIS script for converting the DN pixel values to reflectance for correction of the top-of-atmosphere radiance,style=mystyle,label={lst01}]
#!/bin/sh
r.info -r LC08_L2SP_189056_20131220_20200912_02_T1_SR_B1
gdalinfo LC08_L2SP_189056_20131220_20200912_02_T1_SR_B7.TIF
r.in.gdal LC08_L2SP_189056_20131220_20200912_02_T1_SR_B1.TIF out=L8_2013_01
r.in.gdal LC08_L2SP_189056_20131220_20200912_02_T1_SR_B2.TIF out=L8_2013_02
r.in.gdal LC08_L2SP_189056_20131220_20200912_02_T1_SR_B3.TIF out=L8_2013_03
r.in.gdal LC08_L2SP_189056_20131220_20200912_02_T1_SR_B4.TIF out=L8_2013_04
r.in.gdal LC08_L2SP_189056_20131220_20200912_02_T1_SR_B5.TIF out=L8_2013_05
# repeated likewise until Band 11
g.list rast
g.copy raster=L8_2013_01,lsat8_2013.1
g.copy raster=L8_2013_02,lsat8_2013.2
g.copy raster=L8_2013_03,lsat8_2013.3
g.copy raster=L8_2013_04,lsat8_2013.4
g.copy raster=L8_2013_05,lsat8_2013.5
# repeated likewise until Band 11
i.landsat.toar input=lsat8_2013. output=lsat8_2013_toar. sensor=oli8 \
    method=dos1 date=2013-12-20 sun_elevation=51.85549756 \
    product_date=2013-12-20 gain=HHHLHLHHL
\end{lstlisting}

\subsection[\appendixname~\thesubsection]{GRASS GIS script for computing the ARVI}
\begin{lstlisting}[language=bash,caption=GRASS GIS script for computing the ARVI,style=mystyle,label={lst02}]
g.region raster=lsat8_2013_toar.3 -p
i.vi blue=lsat8_2013_toar.2 red=lsat8_2013_toar.4 nir=lsat8_2013_toar.5 \
       viname=arvi output=lsat8_2013.arvi --overwrite
r.colors lsat8_2013.arvi color=bcyr -e
d.mon wx1
d.rast lsat8_2013.arvi
d.legend raster=lsat8_2013.arvi range=-0.7,0.3 title="ARVI" title_fontsize=14 font=Helvetica fontsize=12 -t -s -b border_color=white thin=12 label_step=0.1 -d
d.out.file output=Nigeria_ARVI_2013 format=jpg --overwrite
\end{lstlisting}

\subsection[\appendixname~\thesubsection]{GRASS GIS script for computing the GARI}
\begin{lstlisting}[language=bash,caption=GRASS GIS script for computing the GARI,style=mystyle,label={lst03}]
g.region raster=lsat8_2013_toar.3 -p
i.vi blue=lsat8_2013_toar.2 green=lsat8_2013_toar.3 red=lsat8_2013_toar.4 \
     nir=lsat8_2013_toar.5 viname=gari output=lsat8_2013.gari --overwrite
r.colors lsat8_2013.gari color=rainbow -e
d.mon wx2
d.rast lsat8_2013.gari
d.legend raster=lsat8_2013.gari range=-0.5,0.7 title="GARI" title_fontsize=14 font="Helvetica" fontsize=12 -t -b bgcolor=white label_step=0.1 border_color=white thin=8 -d
d.out.file output=Nigeria_GARI_2013 format=jpg --overwrite
\end{lstlisting}

\subsection[\appendixname~\thesubsection]{GRASS GIS script for computing the GVI}
\begin{lstlisting}[language=bash,caption=GRASS GIS script for computing the GVI,style=mystyle,label={lst04}]
g.region raster=lsat8_2013_toar.3 -p
i.vi blue=lsat8_2013_toar.2 green=lsat8_2013_toar.3 red=lsat8_2013_toar.4 nir=lsat8_2013_toar.5 band5=lsat8_2013_toar.6 band7=lsat8_2013_toar.7 viname=gvi output=lsat8_2013.gvi --overwrite
r.colors.stddev lsat8_2013.gvi
d.mon wx3
d.rast lsat8_2013.gvi
d.legend raster=lsat8_2013.gvi range=-1,0.7 title="GVI" title_fontsize=14 font="Helvetica" fontsize=12 -t -b bgcolor=white label_step=0.1 border_color=white thin=8 -d
d.out.file output=Nigeria_GVI_2013 format=jpg --overwrite
\end{lstlisting}

\subsection[\appendixname~\thesubsection]{GRASS GIS script for computing the DVI}
\begin{lstlisting}[language=bash,caption=GRASS GIS script for computing the DVI,style=mystyle,label={lst05}]
g.region raster=lsat8_2013_toar.1 -p
i.vi blue=lsat8_2013_toar.2 red=lsat8_2013_toar.4 nir=lsat8_2013_toar.5 viname=dvi output=lsat8_2013.dvi --overwrite
r.colors lsat8_2013.dvi color=bgyr -e
d.mon wx4
d.rast lsat8_2013.dvi
d.legend raster=lsat8_2013.dvi range=-0.1,0.3 title="DVI" title_fontsize=14 font="Helvetica" fontsize=12 -t -b bgcolor=white label_step=0.02 border_color=white thin=8 -d
d.out.file output=Nigeria_DVI_2013 format=jpg --overwrite
\end{lstlisting}

\subsection[\appendixname~\thesubsection]{GRASS GIS script for computing the PVI}
\begin{lstlisting}[language=bash,caption=GRASS GIS script for computing the PVI,style=mystyle,label={lst06}]
g.region raster=lsat8_2013_toar.1 -p
i.vi red=lsat8_2013_toar.4 nir=lsat8_2013_toar.5 viname=pvi output=lsat8_2013.pvi --overwrite
r.colors lsat8_2013.pvi color=bgyr -e
d.mon wx5
d.rast lsat8_2013.pvi
d.legend raster=lsat8_2013.pvi range=-0.1,0.3 title="PVI" title_fontsize=14 font="Helvetica" fontsize=12 -t -b bgcolor=white label_step=0.02 border_color=white thin=8 -d
d.out.file output=Nigeria_PVI_2013 format=jpg --overwrite
\end{lstlisting}

\subsection[\appendixname~\thesubsection]{GRASS GIS script for computing the GEMI}
\begin{lstlisting}[language=bash,caption=GRASS GIS script for computing the GEMI,style=mystyle,label={lst07}]
# Calculation of GEMI: Global Environmental Monitoring Index
g.region raster=lsat8_2013_toar.1 -p
i.vi red=lsat8_2013_toar.4 nir=lsat8_2013_toar.5 viname=gemi output=lsat8_2013.gemi --overwrite
r.colors lsat8_2013.gemi color=roygbiv -e
d.erase
d.rast lsat8_2013.gemi
d.legend raster=lsat8_2013.gemi range=-0.5,1.0 title="GEMI" title_fontsize=14 font="Helvetica" fontsize=12 -t -b bgcolor=white label_step=0.1 border_color=white thin=8 -d
d.out.file output=Nigeria_GEMI_2013 format=jpg --overwrite
\end{lstlisting}

\subsection[\appendixname~\thesubsection]{GRASS GIS script for computing the NDWI}
\begin{lstlisting}[language=bash,caption=GRASS GIS script for computing the NDWI,style=mystyle,label={lst08}]
g.region raster=lsat8_2013_toar.1 -p
i.vi green=lsat8_2013_toar.3 nir=lsat8_2013_toar.5 viname=ndwi output=lsat8_2013.ndwi --overwrite
r.colors lsat8_2013.ndwi color=viridis -e
d.mon wx0
d.rast lsat8_2013.ndwi
d.legend raster=lsat8_2013.ndwi range=-1.0,0.5 title="NDWI" title_fontsize=14 font="Helvetica" fontsize=12 -t -b bgcolor=white label_step=0.1 border_color=white thin=8 -d
d.out.file output=Nigeria_NDWI_2013 format=jpg --overwrite
\end{lstlisting}

\subsection[\appendixname~\thesubsection]{GRASS GIS script for computing the MSAVI2}
\begin{lstlisting}[language=bash,caption=GRASS GIS script for computing the MSAVI2,style=mystyle,label={lst09}]
g.region raster=lsat8_2013_toar.1 -p
i.vi red=lsat8_2013_toar.4 nir=lsat8_2013_toar.5 viname=msavi2 output=lsat8_2013.msavi2 --overwrite
r.colors lsat8_2013.msavi2 color=soilmoisture -e
d.mon wx0
d.rast lsat8_2013.msavi2
d.legend raster=lsat8_2013.msavi2 range=-1.0,0.5 title="MSAVI2" title_fontsize=14 font="Helvetica" fontsize=12 -t -b bgcolor=white label_step=0.1 border_color=white thin=8 -d
d.out.file output=Nigeria_MSAVI2_2013 format=jpg --overwrite
\end{lstlisting}

\subsection[\appendixname~\thesubsection]{GRASS GIS script for computing the IPVI}
\begin{lstlisting}[language=bash,caption=GRASS GIS script for computing the IPVI,style=mystyle,label={lst10}]
g.region raster=lsat8_2013_toar.1 -p
i.vi red=lsat8_2013_toar.4 nir=lsat8_2013_toar.5 viname=ipvi output=lsat8_2013.ipvi --overwrite
r.colors lsat8_2013.ipvi color=haxby -e
d.mon wx0
d.rast lsat8_2013.ipvi
d.legend raster=lsat8_2013.ipvi range=-1.0,1.0 title="IPVI" title_fontsize=14 font="Helvetica" fontsize=12 -t -b bgcolor=white label_step=0.1 border_color=white thin=8 -d
d.out.file output=Nigeria_IPVI_2013 format=jpg --overwrite
\end{lstlisting}

\subsection[\appendixname~\thesubsection]{GRASS GIS script for computing the EVI}
\begin{lstlisting}[language=bash,caption=GRASS GIS script for computing the EVI,style=mystyle,label={lst11}]
g.region raster=lsat8_2013_toar.1 -p
i.vi blue=lsat8_2013_toar.2 red=lsat8_2013_toar.4 nir=lsat8_2013_toar.5 viname=evi output=lsat8_2013.evi --overwrite
r.colors lsat8_2013.evi color=elevation -e
d.mon wx0
d.rast lsat8_2013.evi
d.legend raster=lsat8_2013.evi range=-1.0,1.0 title="EVI" title_fontsize=14 font="Helvetica" fontsize=12 -t -b bgcolor=white label_step=0.1 border_color=white thin=8 -d
d.out.file output=Nigeria_EVI_2013 format=jpg --overwrite
\end{lstlisting}

%%%%%%%%%%%%%%%%%%%%%%%%%%%%%%%%%%%%%%%%%%
%=====================================
% References: external bibliography
%=====================================
\begin{adjustwidth}{-\extralength}{0cm}

\bibliography{Nigeria.bib}
\reftitle{References}
\nocite{*}

\end{adjustwidth}
%%%%%%%%%%%%%%%%%%%%%%%%%%%%%%%%%%%%%%%%%%
\end{document}
